%% PHILOSOPHER ist ein Titlepageprovider der nach und nach die alten provider ersetzen soll

\usetikzlibrary{decorations.text}

\input Konanur.fd
\makeatletter

%% Mathe
\xglobal\definecolor{FacultyMathexColor}{RGB}{086,170,028}
%% Praktische Informatik
\xglobal\definecolor{FacultyPraktischeInformatikxColor}{RGB}{221,101,015}
%% Theoretische Informatik
\xglobal\definecolor{FacultyTheoretischeInformatikxColor}{RGB}{002,133,221}
%% Technische Informatik
\xglobal\definecolor{FacultyTechnischeInformatikxColor}{RGB}{177,2,221}

\newlength{\LILLYxLeftsep}\setlength{\LILLYxLeftsep}{2cm}

%% Upper Text %% Symbol %% WHITE Color %% BACK Color %% Tikz %% Size %% FONTSIZE UPPER
\DeclareDocumentCommand{\LILLYxGENxFACULTY} { m m m O{} O{} O{} O{} O{0.1} O{0.05}} {
    \begin{tikzpicture}[#5]
        \draw [decorate,decoration={text along path,text align=center,
        text={|\footnotesize\normalNumber\color{#3}#7|#1}}]
       (-1.275-#8,0) arc [start angle=180, end angle=00, radius=1.275+#8];
       \draw(-1.5,0) node {{$\bullet$}} ++ (3,0) node {{$\bullet$}};
       \draw [decorate,decoration={text along path,text align=center,
       text={|\footnotesize\normalNumber\color{#3}| ulm university  }}]
      (-1.6,0) arc [start angle=-180, end angle=00, radius=1.6];
       \draw[very thick, color=#3] (0,0) circle (1.25);
       \draw[thick, color=#3] (0,0) circle (1.75);
       \draw[thin, color=#3] circle (1.2);
       \foreach \r in {0,20,...,360} {
           \filldraw[#3] (\r:1.2) circle (0.01) node (\r) {};
       }
       \foreach \r in {0,20,...,360} {
           \foreach \rr in {0,20,...,\r}{
               \draw[very thin] (\r) -- (\rr);
           }
       }
       \node[opacity=0.5] at(0.02+#9,-#9) {\fontsize{#6}{16pt}\selectfont#2};
       \draw (0.02,0) node {\fontsize{#6}{16pt}\selectfont\color{#4}#2};
   \end{tikzpicture}
}

\DeclareDocumentCommand{\LILLYxFACULTYxMATHE}{ O{MudWhite} O{FacultyMathexColor} O{} }{%
\LILLYxGENxFACULTY{Mathematik}{$\pi$}{#1}[#2][#3][80pt][\smallNumber][0.125]
}
\DeclareDocumentCommand{\LILLYxFACULTYxPRAKTISCHEINFORMATIK}{ O{MudWhite} O{FacultyPraktischeInformatikxColor} O{} }{%
    \LILLYxGENxFACULTY{Praktische Informatik}{\hspace*{-0.025em}\bfseries\faCode}{#1}[#2][#3][50pt][\smallNumber][0.125]
}
\DeclareDocumentCommand{\LILLYxFACULTYxTHEORETISCHEINFORMATIK}{ O{MudWhite} O{FacultyTheoretischeInformatikxColor} O{} }{%
    \LILLYxGENxFACULTY{Theoretische Informatik}{\bfseries\hspace*{-0.1em}\faCogs}{#1}[#2][#3][45pt][\smallNumber][0.125]
}
\DeclareDocumentCommand{\LILLYxFACULTYxTECHNISCHEINFORMATIK}{ O{MudWhite} O{FacultyTechnischeInformatikxColor} O{} }{%
    \LILLYxGENxFACULTY{Technische Informatik}{\hspace*{-0.025em}\begin{tikzpicture}\filldraw[rounded corners=4pt] (0,0) rectangle ++(1,1);%
        \foreach \y in {1.117,-0.317} {%
            \foreach \x in {0,0.18,...,0.9} {%%
                \filldraw[rounded corners=0.04] (\x,\y) rectangle ++(0.07,0.2); % TOP/BOT
                \filldraw[rounded corners=0.04] (\y,\x) rectangle ++(0.2,0.07); % LEFT/RIGHT
            }%
        }%
    \end{tikzpicture}}{#1}[#2][#3][50pt][\smallNumber][0.125][0.03]
}

\providecommand{\LILLYxColorxTITLExSETTINGSxGENERAL}{\LILLYxPATHxDATA/Semester/\LILLYxSemester/Definitions/GENERAL.tex}
\providecommand{\LILLYxColorxTITLExSETTINGSxVORLESUNG}{\LILLYxPATHxDATA/Semester/\LILLYxSemester/Definitions/\LILLYxVorlesung}
\ifnum\LILLYxSemester>0
\input{\LILLYxColorxTITLExSETTINGSxGENERAL}\input{\LILLYxColorxTITLExSETTINGSxVORLESUNG}
\fi


%%% Dieser Befehl malt die untere Leiste und setzt die entsprechenden Daten. Hierbei Kann, von der Höhe abgesehen, keine andere Metrik modifiziert werden.
\DeclareDocumentCommand{\LILLYxPHILOSOPHERxBORDERBLOCK}{%
    O{\LILLYxFACULTYxCOLOR} % SignaturFarbe
    O{LightGray}            % Textfarbe
    m%{3}                   % Baseheight
    O{\LILLYxFlavourText}   % Textblahblah
    O{\LILLYxFACULTY}       % Zu setzendes Symbol
}{%
\begin{tikzpicture}%
    %\node[above right] at (\LILLYxLeftsep,3.05) {\centering\fontsize{3pt}{2pt}\selectfont\centering Fassung vom \heute{} um \currenttime Uhr\par};
    \draw[thick,#1] (0,#3+0.1) -- ++ (0.6\paperwidth-0.1,0) -- ++(0.1\paperwidth,1) -- ++(0.1\paperwidth+0.1,0) --  ++(0.1\paperwidth,-1) -- ++(0.1\paperwidth,0) ;
    \draw[thin,#1] (0,#3+0.175) -- ++ (0.6\paperwidth-0.175,0) -- ++(0.1\paperwidth,1) -- ++(0.1\paperwidth+0.175,0) --  ++(0.1\paperwidth,-1) -- ++(0.1\paperwidth,0) ;
    \filldraw[#1] (0,0) -- ++(\paperwidth,0) -- ++(0,#3) --  ++(-0.1\paperwidth,0) -- ++(-0.1\paperwidth,1) -- ++(-0.1\paperwidth,0) --  ++(-0.1\paperwidth,-1) -- ++(-0.6\paperwidth,0) -- cycle;%
    \begin{scope}[every node/.style={color=#2}]%
        \node at(0.75\paperwidth,2) {#5}; %\resizebox{!}{#3cm}{
        \node[right] at(\LILLYxLeftsep,1.5) {\parbox{0.45\linewidth}{\footnotesize{#4\par}}};
        \node[above left] at(\paperwidth,0.02) {\begingroup\hypersetup{urlcolor=#2}\href{https://github.com/EagleoutIce/LILLY}{\usefont{U}{Konanur}{xl}{n}LILLY}\hypersetup{urlcolor=\LILLYxColorxLINKSxUrlColor}\endgroup};%
        \node[above left] at(\paperwidth,0.52) {\tiny \textsf{Ver. \LILLYxVERSION}};%
    \end{scope}%
\end{tikzpicture}%
}

\gdef\LILLYxPHILOSOPHERxINIT{\newgeometry{margin = 0pt}}

%% This variant prints an highlighted Bonus and include \LILLY@n/a number if wanted as an extra number to seperate Versions
\DeclareDocumentCommand{\LILLYxTITLExBONUS}{%
    O{\LILLYxFACULTYxCOLOR} % SignaturFarbe
    O{MudWhite}             % Textfarbe
    m                       % Main Matter => Bonus text upper left
    O{\LILLYxFlavourText}   % Textblahblah
    O{\LILLYxFACULTY}       % Zu setzendes Symbol
    O{0.55\paperwidth}      % Scaling
    O{\LILLYxPATHxDATA/Semester/\LILLYxSemester/Graphics/titleimage\LILLYxVorlesung.pdf}
                            % Titel Symbol
}{%
\begingroup\def\old@textwidth{0.75\paperwidth}%
\clearpage\begin{titlepage}
    \LILLYxPHILOSOPHERxINIT%
    {\par\vspace*{6em}\par}%
    {% Upper title Block. It can hold the denoted 'Main Matter'
\begin{tikzpicture}
    % Halo (3 confirmed hihi):
    \fill[#1,opacity=0.5] (0,0) |- ++(#6+0.4,1.25) -- ++(-0.77,-1.35) -- cycle;
    % Texture:
    \filldraw[#1] (0,0) |- ++(#6,1.25) -- ++(-0.75,-1.25) -- cycle;
    \node[right,#2] at (\LILLYxLeftsep,0.5) {\Huge\bfseries#3};
    \begin{pgfinterruptboundingbox}
    \node[left] at(\paperwidth-1.5cm,0.25) {\includegraphics[height=11em]{#7}};
    \end{pgfinterruptboundingbox}
\end{tikzpicture}
    } {\par\vspace*{9em}\par}%
    {\centering% Title:
    {\bfseries\usefont{T1}{qzc}{m}{it}\fontsize{40pt}{16pt}\selectfont\FULLTITLE}\par%
    } {\par\vspace*{3em}\par}%
    {\centering% Metadatame - hihih:
    {\scshape\tiny{#3 von}\\\large\AUTHOR\par}
	\vspace{0.5cm}
    {\itshape\small Version vom:\\ \heute\par}
    }\hbox{}\vfill\vspace{4em}
    %{\listofTOPICS}
    {\centering\parbox{\old@textwidth}{\let\textwidth\old@textwidth{}\POLITEINTRO}\par}
    \hbox{}\vfill%
    \LILLYxPHILOSOPHERxBORDERBLOCK[#1][#2]{3}[#4][#5] %% height
\end{titlepage}%
\endgroup\restoregeometry%
}

\DeclareDocumentCommand{\LILLYxTITLExRAW}{ o }{
    \begingroup%
    \clearpage\begin{titlepage}
        \LILLYxPHILOSOPHERxINIT%
        {\centering\par\vspace*{7em}\par
            {\raggedright\hspace*{\LILLYxLeftsep}\parbox{\linewidth-\LILLYxLeftsep}{\bfseries\usefont{T1}{qzc}{m}{it}\fontsize{40pt}{16pt}\selectfont\FULLTITLE}\par}
            \par\vspace*{11em}\par
            \includegraphics[height=20em]{#1}
            \par\vfill\vspace{2em}\par
            {\raggedright\hspace*{\LILLYxLeftsep}\parbox{\linewidth-\LILLYxLeftsep}{{\scshape \tiny{\LILLY@Typ{} von}\\\large\AUTHOR\\[-0.35\baselineskip]{\normalfont\LILLYxlstTypeWriter\tiny\AUTHORMAIL}\par}
            \vspace{0.5cm}
            {\itshape\small Version vom: \\ \heute\par}}}
        }
        \hbox{}\vfill
        \LILLYxPHILOSOPHERxBORDERBLOCK{3} %% height
    \end{titlepage}\endgroup\restoregeometry%
}

%% Pfad zu den Einstellungen %% Pfad zur Grafik
\DeclareDocumentCommand{\LILLYxTITLE}{ O{\LILLYxPATHxDATA/Semester/\LILLYxSemester/Graphics/titleimage\LILLYxVorlesung.pdf} } {%
    \ifnum\LILLYxSemester>0\LILLYxTITLExRAW[#1]\fi
}